\documentclass[11pt]{article}

\usepackage[a4paper,margin=1in]{geometry}
\usepackage{amsmath,amssymb,amsthm,mathtools}
\usepackage{enumitem}
\usepackage[hidelinks]{hyperref}

\newtheorem{theorem}{Theorem}[section]
\newtheorem{proposition}[theorem]{Proposition}
\newtheorem{lemma}[theorem]{Lemma}
\newtheorem{corollary}[theorem]{Corollary}
\theoremstyle{definition}
\newtheorem{remark}[theorem]{Remark}
\newtheorem{problem}[theorem]{Problem}

\newcommand{\R}{\mathbb R}
\newcommand{\Om}{\Omega}
\newcommand{\Sph}{\mathbb S^{n-1}}
\newcommand{\Drr}{\mathcal D_{\rm rr}}
\newcommand{\Rrad}{\mathcal R_{\rm rad}}
\newcommand{\dT}{\delta_T}
\newcommand{\Fr}{\mathcal A}
\newcommand{\HH}{\mathcal H}
\newcommand{\X}{\mathcal X}
\newcommand{\eps}{\varepsilon}
\newcommand{\norm}[1]{\left\|#1\right\|}
\newcommand{\abs}[1]{\left|#1\right|}

\title{Debugging the Continuous-Rearrangement Proof\\
       \large What survives, what fails, and the quantitative route that remains}
\author{}
\date{May 2026}

\begin{document}
\maketitle

\begin{abstract}
This note audits \texttt{faber\_krahn\_rearrangement\_analysis.tex}.
The conclusion is not that continuous rearrangement is useless; the
conclusion is that the current Brock/Steiner schedule proof asks the
wrong object to do the sharp endpoint job.  The finite-time dense-schedule
claim is false, the stated one-dimensional quantitative
Polya--Szego theorem is false as written, the multidimensional Fubini
transition ignores transverse derivatives, and the centre-drift lemma is
false.  The repair is to replace the one-shot directional story by the
Plan~2 continuous \emph{level-set} rearrangement: Talenti's level-set
identity gives the exact budget
\[
   \dT(\Omega)\simeq \int (D_H+D_I),
\]
where \(D_I\) controls distance of many levels from balls and \(D_H\)
controls the kinetic cost of moving between levels.  A weighted metric
trace then gives a good near-boundary level; the boundary-layer transfer
returns the estimate to \(\Omega\).  This is the route in the directory
with the correct \(H^{1/2}\)/quadratic scale.  It is fully quantitative
once the explicitly isolated remaining lemmas are proved with constants.
\end{abstract}

\tableofcontents

\section{The target and the non-negotiable scale}

Normalize \(|\Omega|=\omega_n\), and write \(u=u_\Omega\).  The sharp
Saint--Venant/Faber--Krahn target is
\begin{equation}\label{eq:target}
   \Fr(\Omega)^2 \le C(n)\,\dT(\Omega).
\end{equation}
The rearrangement bookkeeping begins with the exact splitting
\begin{equation}\label{eq:split}
   \dT(\Omega)
   =
   \underbrace{\left(\int |\nabla u|^2-\int |\nabla u^*|^2\right)}_{\Drr(\Omega)}
   +
   \underbrace{\left(T(B)-J(u^*)\right)}_{\Rrad(\Omega)}.
\end{equation}
This identity is sound.  The radial profile term is harmless: by strict
concavity of the torsion functional on the ball,
\[
   \Rrad(\Omega)=\int_B |\nabla(u^*-u_B)|^2.
\]
The question is how to convert \(\Drr\), or a refined form of it, into
geometric asymmetry at the sharp quadratic scale.

\section{Where the current proof breaks}

\subsection{Finite-time dense schedules do not see all directions}

The theorem called ``Lower Bound for Dense Schedules'' in the current
rearrangement-analysis note is false for finite time.  If the step size
is fixed and \(T<\Delta_0\), the flow uses only one Steiner direction.
Any non-radial function already symmetric in that direction has zero
dissipation on \([0,T]\) but positive distance from the radial cone.
More generally, a finite interval samples only finitely many directions;
it never samples a dense subset of \(\Sph\).

Even if the statement is weakened to an infinite-time convergence
statement, the proof is not quantitative.  The compactness argument would
produce at best a modulus depending on a compact class, the schedule, and
the topology.  It does not produce the computable quadratic constant
needed for \eqref{eq:target}.  The missing theorem would have to be a
genuine angular spectral-gap estimate for averaged instantaneous
dissipation, not compactness.

\subsection{The stated one-dimensional theorem is false}

The note states a scale-free estimate
\begin{equation}\label{eq:false-1d}
   \int |g'|^2-\int |(g^*)'|^2
   \ge c_0 \inf_z \norm{g-g^*(\cdot-z)}_{L^2(\R)}^2
\end{equation}
for all compactly supported \(H^1\) functions on the line.  This cannot
hold with a universal \(c_0\).  If \(g_L(x)=g(x/L)\), then the left side
scales like \(L^{-1}\), while the right side scales like \(L\).  The
ratio is \(L^{-2}\).

There is a second, structural error: non-symmetric does not imply that
the level sets \(\{g>t\}\) have multiple intervals on a positive set of
heights.  A single asymmetric tent function has one interval at every
height, but is not symmetric decreasing up to translation.  The
multi-interval Cauchy--Schwarz estimate is correct but incomplete.

\subsection{The missing one-dimensional term is centre drift}

For one regular interval \(\{g>t\}=(a(t),b(t))\), set
\[
   p(t)=\frac{1}{g'(a(t))}>0,\qquad
   q(t)=\frac{-1}{g'(b(t))}>0.
\]
Then \(a'=p\), \(b'=-q\), the length \(\ell=b-a\) satisfies
\(-\ell'=p+q\), and the centre \(c=(a+b)/2\) satisfies
\[
   c'=\frac{p-q}{2}.
\]
The exact one-interval Polya--Szego level loss is
\begin{equation}\label{eq:one-interval-loss}
   \frac1p+\frac1q-\frac{4}{p+q}
   =
   \frac{(p-q)^2}{pq(p+q)}
   =
   \frac{4|c'|^2}{pq(p+q)}.
\end{equation}
Thus a one-interval asymmetric slice is not an obstruction by itself; it
is controlled through variation of the level-centre \(c(t)\), provided
the weights are non-degenerate on the height interval under study.  This
is precisely why a correct proof needs profile-gap and bad-density
control.  The current proof underclaims here by treating multiplicity as
the only source of dissipation, and overclaims by pretending the
multiplicity term controls all asymmetry.

\subsection{The multidimensional Fubini step is not exact}

For a fixed direction \(\nu\), slicing gives an identity for the
\(\partial_\nu\)-part of the gradient, not for the full Dirichlet energy:
\[
   \int_{\R^n}|\partial_\nu u|^2-\int_{\R^n}|\partial_\nu u^{*\nu}|^2
   =
   \int_{\nu^\perp}
   \left(\int |(u_{y,\nu})'|^2-\int |((u_{y,\nu})^*)'|^2\right)\,dy.
\]
The full difference
\[
   \int|\nabla u|^2-\int|\nabla u^{*\nu}|^2
\]
also involves transverse derivatives.  Brock's theorem controls the full
Dirichlet energy under continuous Steiner symmetrization, but it is not a
plain Fubini identity.  Any proof that averages one-dimensional slice
losses must either use Brock's full formula or switch to a level-set
identity where all terms are accounted for.

\subsection{The centre-drift quarantine lemma is false}

The asserted estimate
\[
   \int_{\nu^\perp} |\bar s_\nu(y)-x_0\cdot\nu|^2M_\nu(y)\,dy
   \lesssim
   \int|\nabla u|^2-\int|\nabla u^{*\nu}|^2
\]
is false.  Take a function whose slices in the \(s\)-direction are
translated copies of a fixed symmetric profile, with the translation
depending on \(y\), and with negligible \(y\)-transition cost in separated
components.  The one-direction Steiner gap can be zero on each slice,
while the variance of slice centres is positive.  The proof also uses an
infinite factor \(\int_\R s^2\,ds\), an unavailable Hopf lower bound for
all slice masses, and a Poincare estimate for \(u-u^{*\nu}\) that is not
true in the stated form.

\section{The repaired quantitative object}

The route in the directory that has the right scale is not the finite
directional Brock schedule.  It is the Talenti/Plan~2 continuous
level-set rearrangement.  Let
\[
   E_t:=\{u>t\},\qquad m(t):=|E_t|.
\]
For a.e. regular level define
\[
   D_H(t)
   :=
   \left(\int_{\partial^*E_t}\frac{1}{|\nabla u|}\,d\HH^{n-1}\right)
   \left(\int_{\partial^*E_t}|\nabla u|\,d\HH^{n-1}\right)
   -P(E_t)^2,
\]
\[
   D_I(t):=P(E_t)^2-n^2\omega_n^{2/n}m(t)^{2-2/n}.
\]
The first term is the Holder/Cauchy--Schwarz loss; the second is the
isoperimetric loss of the level set.  The proved level-set identity in
the Plan~2 material is
\begin{equation}\label{eq:level-budget}
   \frac{2\dT(\Omega)}{\omega_n}
   =
   \frac{1}{\omega_n}\int_0^{\|u\|_\infty}
   \frac{D_H(t)+D_I(t)}
        {n^2\omega_n^{2/n}m(t)^{1-2/n}}\,dt.
\end{equation}
This is the quantitative replacement for the informal statement
``the rearrangement flow dissipates energy''.  It is exact, levelwise,
and separates the two mechanisms that the broken proof conflated:
\begin{itemize}[leftmargin=2em]
\item \(D_I\) says that many levels are close to balls;
\item \(D_H\) says that the level flow has controlled normal-speed
oscillation, hence controlled kinetic motion between nearby levels.
\end{itemize}

\section{The debugged route to the sharp theorem}

Parametrize levels by volume radius:
\[
   |E_\rho|=\omega_n\rho^n,\qquad
   F_\rho:=\rho^{-1}E_\rho.
\]
Work in the quotient metric space
\[
   d_\X([A],[C])=\inf_{z\in\R^n}|A\Delta(C+z)|.
\]
The correct sharp route is the following.

\begin{enumerate}[label=\textbf{Step \arabic*.},leftmargin=3em]
\item \textbf{Exact budget.}
Use \eqref{eq:level-budget} to obtain, on any fixed bulk window
\(\rho\in[\rho_*,1]\),
\[
   \int_{\rho_*}^{\rho_\delta}
   (D_H(t(\rho))+D_I(t(\rho)))\,d\mu(\rho)
   \le C(n,\rho_*)\,\dT(\Omega),
\]
where \(d\mu(\rho)=(-t'(\rho))\,d\rho\).

\item \textbf{Distance on many levels from \(D_I\).}
On levels with small \(D_I\), a same-centre strong quantitative
isoperimetric estimate gives
\[
   d_\X(F_\rho,B_1)^2\le C\,D_I(t(\rho)).
\]
This is where the strong Fusco--Julin type input belongs.

\item \textbf{Kinetic control from \(D_H\).}
For a Borel centre field \(z_\rho\), the desired metric first variation is
\[
   |\dot F_\rho|_\X
   \le
   \rho^{-n}\inf_{a\in\R^n}
   \int_{\partial^*E_\rho}
   |V_\rho-H_{z_\rho,\rho}-a\cdot\nu_\rho|\,d\HH^{n-1},
\]
where
\[
   V_\rho=\frac{-t'(\rho)}{|\nabla u|},\qquad
   H_{z_\rho,\rho}(x)=\frac{(x-z_\rho)\cdot\nu_\rho(x)}{\rho}.
\]
The self-truncated velocity lemma in the directory shows that the
low-gradient part is paid for by \(D_H\); no Hopf lower bound is needed.
Together with the same-centre homothetic trace this gives
\[
   |\dot F_\rho|_\X^2\le C(D_H(t(\rho))+D_I(t(\rho)))
\]
on good levels.

\item \textbf{Weighted trace.}
The profile-gap lemma gives a density lower bound for \(d\mu\) outside a
bad set of Lebesgue measure \(O(\dT)\) near
\[
   \rho_\delta=1-\kappa\sqrt{\dT(\Omega)}.
\]
Combining Markov on the distance term with Cauchy--Schwarz on the kinetic
term over an order-one window gives
\[
   d_\X(F_{\rho_\delta},B_1)^2\le C\,\dT(\Omega),
\]
or the good-endpoint variant gives the same estimate at
\(\widehat\rho\in[\rho_\delta-C\dT,\rho_\delta]\).

\item \textbf{Boundary-layer transfer.}
Since \(E_{\widehat\rho}\subset\Omega\) and
\[
   |\Omega\setminus E_{\widehat\rho}|
   =
   \omega_n(1-\widehat\rho^n)
   \le C\sqrt{\dT(\Omega)},
\]
the elementary symmetric-difference transfer gives
\[
   \Fr(\Omega)
   \le
   \Fr(E_{\widehat\rho})
   +C\sqrt{\dT(\Omega)}.
\]
The metric trace gives
\[
   \Fr(E_{\widehat\rho})^2\le C\,\dT(\Omega).
\]
Squaring the transfer inequality proves the sharp estimate
\[
   \Fr(\Omega)^2\le C\,\dT(\Omega).
\]
\end{enumerate}

This route is a continuous rearrangement proof in the right sense: the
``time'' is the level/volume radius, the exact energy budget is
\eqref{eq:level-budget}, and the proof is an action-plus-trace argument.
It is not the finite-time dense-direction Brock proof.

\section{What remains to make it fully quantitative}

The remaining tasks are now clean.  If each item below is proved with
constants, then all constants in the final theorem are computable by
composition.

\begin{problem}[Same-centre strong isoperimetric input]\label{prob:same-centre}
For every finite-perimeter set \(E\), prove or cite in exactly this
form that there is one centre \(z_E\) such that
\[
   \left(\frac{|E\Delta B_E(z_E)|}{|E|}\right)^2
   +
   \frac{1}{P(B_E)}
   \int_{\partial^*E}
   \left|\nu_E-\frac{x-z_E}{|x-z_E|}\right|^2\,d\HH^{n-1}
   \le C(n)\,\delta_{\rm iso}(E).
\]
The strong theorem of Fusco--Julin controls Fraenkel asymmetry and
boundary oscillation, but the route needs the \emph{same} centre in both
terms.  If the cited statement is only separately minimized, a centre
comparison lemma must be added.
\end{problem}

\begin{problem}[Finite-perimeter metric upper gradient]\label{prob:UG}
Upgrade the smooth first-variation formula for
\(\rho\mapsto [\rho^{-1}E_\rho]\) to finite-perimeter level sets:
\[
   |\dot F_\rho|_\X
   \le
   \rho^{-n}\inf_{a\in\R^n}
   \int_{\partial^*E_\rho}
   |V_\rho-H_{z_\rho,\rho}-a\cdot\nu_\rho|\,d\HH^{n-1}
\]
for a.e. \(\rho\).  Lower semicontinuity is the wrong direction; the
missing statement is a genuine limsup recovery theorem for the boundary
action, or a direct finite-perimeter deformation estimate.
\end{problem}

\begin{problem}[Bad-level kinetic control]\label{prob:bad}
Control the kinetic action on the bad isoperimetric or bad-density
levels, or avoid transporting through those levels.  The good-endpoint
variant shows that the final level may be chosen in
\(G_I\cap G_\tau\), but it does not remove the need for kinetic control
when propagating the Markov level to that endpoint.  A measure estimate
for the bad-density set is not enough; one needs the right \(L^2(d\rho)\)
action scale.
\end{problem}

\begin{problem}[Globalization]\label{prob:global}
After the bounded result is proved, globalize it by a quantitative
bounded reduction.  The existing final assembly uses the constructive BDV
truncation lemmas and a suboptimal rearrangement bound, not BDV's final
sharp stability theorem.  If the aim is to avoid every BDV input, those
truncation lemmas must be reproved independently with explicit constants.
\end{problem}

\section{Verdict}

The proof in \texttt{faber\_krahn\_rearrangement\_analysis.tex} cannot be
patched by minor edits.  Its main theorem-level claims are false as
stated.  The useful idea inside it is the dissipation picture, but the
fully quantitative implementation must be the Plan~2 level-set action
route, not the finite dense-direction schedule.

The sharp theorem is therefore reduced to Problems
\ref{prob:same-centre}--\ref{prob:bad} in the bounded setting, followed
by Problem~\ref{prob:global} if one wants the global theorem.  This is a
genuine quantitative route: no compactness modulus is needed in the final
trace argument, and all constants are computable once the named inputs
are proved quantitatively.

\begin{thebibliography}{9}

\bibitem{FuscoJulin}
N.~Fusco and V.~Julin,
\emph{A strong form of the quantitative isoperimetric inequality},
Calc. Var. Partial Differential Equations 50 (2014), 925--937;
arXiv:1111.4866.

\bibitem{Plan2LSI}
\texttt{Manuscript/plan2\_levelset\_identity.tex}.
Exact level-set identity, profile-gap estimate, and superlevel transfer.

\bibitem{Plan2WMT}
\texttt{Plan 2/WIP/WIP\_WeightedMetricTrace.tex}.
Conditional integrated action and weighted metric trace.

\bibitem{Plan2MF}
\texttt{Plan 2/WIP/WIP\_MetricFramework.tex}.
Metric first variation and the finite-perimeter upper-gradient gap.

\bibitem{Plan2TR}
\texttt{Plan 2/WIP/WIP\_TrimmedVelocityRepair.tex}.
Self-truncated velocity defect controlled by \(D_H\).

\end{thebibliography}

\end{document}
